\documentclass[10pt,a4paper]{amsart}
\usepackage[margin=19mm]{geometry}
\usepackage{amsmath,amssymb,mathtools}
\usepackage[scr=rsfs,cal=euler]{mathalfa}
\usepackage{xfrac}
\usepackage[unicode,hidelinks,bookmarks=false]{hyperref}
\newcommand{\ie}{i.e.\ }
\newcommand{\cf}{cf.\ }
\newcommand{\resp}{respectively\ }
\newcommand{\Subsection}{\subsection}
\providecommand{\nobreakdash}{\nobreak\hbox{-}}
\setlength{\emergencystretch}{2em}
\multlinegap0pt
\usepackage{bm}
\usepackage{bbm}
\usepackage{dsfont}
\usepackage{xspace}
\usepackage{cancel}
\usepackage{mathtools}
\usepackage{amsthm}
\makeatletter
\@ifundefined{theorem}{\newtheorem{theorem}{Theorem}}{}
\@ifundefined{proposition}{\newtheorem{proposition}[theorem]{Proposition}}{}
\makeatother
\title[The Minimum Cardinality of a Dependent Finite Gabor System I]{The Minimum Cardinality of a Dependent Finite Gabor System Is Four}
\author{Xinan Dai, Wenhao Deng, Yingdong Shi, Tailin Wu, Yuchen Yang}
\date{Source: arXiv:2608.08190v1 (CC BY 4.0)}
\begin{document}
\maketitle

\noindent\emph{Source and licence.} The Minimum Cardinality of a Dependent Finite Gabor System Is Four. Version arXiv:2608.08190v1 (CC BY 4.0), distributed under the Creative Commons Attribution 4.0 International licence, as recorded in the arXiv licence metadata for this exact version and shown on the abstract page https://arxiv.org/abs/2608.08190v1. Full rights notice: licenses/arxiv-2608.08190-CC-BY-4.0.txt.\par\smallskip

\noindent\textit{Editorial scope.} Counted unit: the Proposition whose source label is \texttt{prop:explicit-square} in the article (source0.tex lines 543--557, source label \texttt{prop:explicit-square}), delivered together with the article's own complete original proof of it (source0.tex lines 559--695, one \texttt{proof} environment of 137 lines). No mathematics is written, repaired or re-derived: statement and proof are the article's own text line for line. Required context retained uncounted: eq:explicit-tau (lines 158--162); eq:cert-numbers (lines 478--484); eq:Ptau (lines 525--528); eq:pert-k (lines 532--537). Every definition, display and label of those lines is retained unchanged. Omitted from this package and not redistributed: 1--157, 163--477, 485--524, 529--531, 538--542, 558--558, 696--1029. Nothing inside a retained excerpt is omitted or altered: every retained body is delivered exactly as the article's lines are written, so no omission receipt of any kind accompanies this package. Citations: the retained excerpts carry no citation command at all, so no bibliography entry of the article is transcribed and no reference is redistributed. Reference adaptation: none was needed and none was made; every reference inside the delivered unit resolves inside it, so no unresolved reference or citation is produced. Printed slips kept literal: no printed word, symbol, hypothesis or number of the counted statement or of its proof was altered. Layout and class: the article's own class is \texttt{amsart} with packages including fontenc, amsmath, amssymb, amsthm, mathtools, microtype, hyperref, geometry; the wrapper is the lane's pinned \texttt{amsart} editorial wrapper at 10pt on A4, which restates the theorem-like environments the retained text needs and adds the article's own macro definitions that the retained text needs (none), with extra wrapper packages (bm, bbm, dsfont, xspace, cancel, mathtools). The delivered numbering is the unit's own. Assets: no retained span contains a figure environment, a graphics-inclusion command, a PDF-page inclusion, a table or a TikZ picture; no image file is copied into this package, the package declares no asset dependency and no image file is redistributed with it. Licence: version arXiv:2608.08190v1 of this article is distributed under the Creative Commons Attribution 4.0 International licence, as recorded in the arXiv licence metadata for this exact version and shown on the abstract page https://arxiv.org/abs/2608.08190v1. A full rights notice is delivered at licenses/arxiv-2608.08190-CC-BY-4.0.txt. Characters: every retained excerpt is pure ASCII, so the delivered engine, XeLaTeX with its default Latin Modern text font, reports no missing character.
\begin{equation}\label{eq:explicit-tau}
 \alpha=\frac13+10^{-12}\sqrt2,
 \qquad
 \beta=\frac13+10^{-12}\sqrt3.
\end{equation}

\begin{equation}\label{eq:cert-numbers}
 d\ge 15.4044994528652104323,
 \quad
 b\le0.2499996482177990124,
 \quad
 c\le0.7725133354867000155,
\end{equation}

\begin{equation}\label{eq:Ptau}
 P_\tau(z)=A(z)A(z-\tau)A(z-2\tau)X_0Z_0:
 E_{z-3\Delta}\longrightarrow E_z.
\end{equation}

\begin{align}\label{eq:pert-k}
 k_{11}^\tau&=\langle h(w),J_\tau(z)h(z)\rangle,&
 k_{12}^\tau&=\langle h(w),J_\tau(z)n(z)\rangle,\notag\\
 k_{21}^\tau&=\langle n(w),J_\tau(z)h(z)\rangle,&
 k_{22}^\tau&=\langle n(w),J_\tau(z)n(z)\rangle.
\end{align}

\begin{proposition}[an explicit continuation square]\label{prop:explicit-square}
If
\begin{equation}\label{eq:explicit-square}
 \|\tau-\theta\|_\infty\le 2\cdot10^{-12},
\end{equation}
then the inverse three-step projective return maps the closed graph disk
$|\xi|\le1/4$ strictly into itself and is a strict contraction there.  It has
a unique invariant graph $L_\tau$.  The line bundle $L_\tau$ is smooth,
topologically trivial, depends continuously on $\tau$, and satisfies both
\begin{align}
 D_\tau(z)L_\tau(T_\tau^3z)&=L_\tau(z),                 \label{eq:three-inv}\\
 C_\tau(z)L_\tau(T_\tau z)&=L_\tau(z).                 \label{eq:one-inv}
\end{align}
In particular, the parameter \eqref{eq:explicit-tau} is admissible.
\end{proposition}

\begin{proof}
We make the perturbation estimate deliberately coarse.  On the unit torus,
\begin{equation}\label{eq:A-norms}
 \|A\|\le2,
 \qquad
 \|\partial_xA\|,\ \|\partial_\omega A\|\le\frac\pi2.
\end{equation}
Only the second and third factors of \eqref{eq:Ptau} move with $\tau$.
A telescoping product estimate therefore gives
\begin{equation}\label{eq:P-pert}
 \|P_\tau(z)-P_\theta(z)\|
 \le12\pi\delta<40\delta.
\end{equation}
For $2\times2$ matrices, adjugation is linear in the entries and preserves
the spectral norm up to the fixed unitary transpose identification.  Hence
\begin{equation}\label{eq:J-bounds}
 \|J_\tau-J_\theta\|<40\delta,
 \qquad
 \|J_\tau\|\le8.
\end{equation}
The explicit unit section $h_0$ satisfies
\[
 \|h_0\|=1,
 \qquad
 \|\partial_xh_0\|\le\frac{3\pi}{2},
 \qquad
 \|\partial_\omega h_0\|\le\pi.
\]
Differentiating the rational product gives
$\|\partial_xJ_\theta\|,\|\partial_\omega J_\theta\|\le6\pi$.  Since
$h=J_\theta h_0$, it follows that
\begin{equation}\label{eq:h-bounds}
 \|h\|=\|n\|\le8,
 \qquad
 \|\partial_xh\|,\ \|\partial_\omega h\|<60,
\end{equation}
and the same derivative bound holds for $n$.  Since
$\|w-z\|_\infty\le3\delta$,
\begin{equation}\label{eq:frame-shift}
 \|h(w)-h(z)\|,\ \|n(w)-n(z)\|<400\delta.
\end{equation}
Combining \eqref{eq:J-bounds}--\eqref{eq:frame-shift} in
\eqref{eq:pert-k} gives, uniformly in $z$,
\begin{equation}\label{eq:k-perturbation}
 |k_{ij}^\tau-k_{ij}^\theta|<30000\delta
 \qquad(1\le i,j\le2).
\end{equation}
Under \eqref{eq:explicit-square} we may therefore set
\[
 \eta:=6\cdot10^{-8}
\]
and use $|k_{ij}^\tau-k_{ij}^\theta|\le\eta$.

Let $r_0=1/4$, and denote the three rational certificate bounds in
\eqref{eq:cert-numbers} by
\[
 D=15.4044994528652104323,
 \quad B=0.2499996482177990124,
 \quad C=0.7725133354867000155.
\]
If
\[
 d_\tau=|k_{11}^\tau|-r_0|k_{12}^\tau|,
 \qquad
 N_\tau=|k_{21}^\tau|+r_0|k_{22}^\tau|,
\]
then
\begin{align*}
 d_\tau&\ge d_\theta-(1+r_0)\eta,\\
 N_\tau&\le N_\theta+(1+r_0)\eta.
\end{align*}
At the rational parameter,
$N_\theta\le B d_\theta$ and $d_\theta\ge D$.  Hence
\begin{align}\label{eq:image-margin}
 r_0d_\tau-N_\tau
 &\ge (r_0-B)D-(1+r_0)^2\eta\\
 &>5.3\cdot10^{-6}>0.\notag
\end{align}
Thus the perturbed projective map still sends the closed radius-$1/4$ graph
disk strictly into itself.

It remains to retain contraction.  From \eqref{eq:h-bounds} and
\eqref{eq:J-bounds}, the rational coefficients satisfy
$|k_{ij}^\theta|\le512$.  Consequently
\begin{equation}\label{eq:detK-pert}
 \left|
 (k_{11}^\tau k_{22}^\tau-k_{12}^\tau k_{21}^\tau)
 -(k_{11}^\theta k_{22}^\theta-k_{12}^\theta k_{21}^\theta)
 \right|
 \le2048\eta+2\eta^2.
\end{equation}
Moreover
$d_\tau\ge D-(1+r_0)\eta>15.4044$.  Using
$|\det(k_{ij}^\theta)|\le C d_\theta^2$ and
$d_\theta\le d_\tau+(1+r_0)\eta$, equations
\eqref{eq:detK-pert} and \eqref{eq:cert-numbers} give
\begin{equation}\label{eq:pert-contraction}
 \sup_{z,\,|\xi|\le1/4}|\partial_\xi\Phi_{\tau,z}(\xi)|<0.773<1.
\end{equation}
This proves a uniform contraction on the whole parameter square.

The contraction mapping theorem now gives a unique continuous invariant graph
$L_\tau$, and the usual fixed-point estimate gives continuous dependence on
$\tau$.  Projection of the graph onto $H$ is a bundle isomorphism, so
$L_\tau$ is trivial.  Smoothness in $z$ follows directly by differentiating
the graph iteration: at derivative order $j$, the only term containing the
$j$th derivative of the previous iterate is multiplied by
$\partial_\xi\Phi_{\tau,z}$, whose norm is uniformly $<0.773$; all remaining
terms contain only lower derivatives.  Induction in $j$ therefore gives
uniform convergence of every derivative and hence $L_\tau\in C^\infty$.

For a two-dimensional invertible cocycle, uniform contraction of the inverse
projective graph transform on an invariant graph is the standard graph-transform
criterion for a dominated splitting; the invariant graph is its unique stable
line.  Hence the fixed graph above is the stable line of a dominated splitting
for the three-step cocycle.  The cocycle identity
\begin{equation}\label{eq:cyclic}
 D_\tau(z)C_\tau(T_\tau^3z)=C_\tau(z)D_\tau(T_\tau z)
\end{equation}
shows that
\[
 \widetilde L_\tau(z):=C_\tau(z)L_\tau(T_\tau z)
\]
is another stable line for $D_\tau$.  Indeed, iterating
\eqref{eq:cyclic} conjugates the three-step dynamics on
$L_\tau(T_\tau z)$ to that on $\widetilde L_\tau(z)$ by the bounded invertible
map $C_\tau$.  The stable line of a dominated splitting is unique, so
$\widetilde L_\tau=L_\tau$.  This is \eqref{eq:one-inv}.

Finally, for \eqref{eq:explicit-tau},
\[
 \|\tau-\theta\|_\infty
 =10^{-12}\max\{\sqrt2,\sqrt3\}
 <2\cdot10^{-12},
\]
which proves admissibility of the announced parameter.
\end{proof}

\end{document}
